\documentclass[12pt,a4paper]{article}

% ------------------------------------------------
% Pacotes
% ------------------------------------------------
\usepackage[utf8]{inputenc}
\usepackage[T1]{fontenc}
\usepackage[brazil]{babel}
\usepackage{graphicx}
\usepackage[table]{xcolor}
\usepackage{array}
\usepackage{tabularx}
\usepackage{geometry}
\usepackage{fancyhdr} 
\usepackage{mathptmx} % Times New Roman 
\usepackage{enumitem}
\usepackage{amsfonts}
\usepackage{amsmath}
\usepackage{amssymb}
\usepackage{tcolorbox}

% ------------------------------------------------ 
% Configuração de página
% ------------------------------------------------
\geometry{
    a4paper,
    left=2cm,
    right=2cm,
    top=2cm,
    bottom=2cm,
    headheight=15pt
}

\pagestyle{fancy}
\fancyhf{}
\renewcommand{\headrulewidth}{0pt}
\setlength{\parindent}{0pt}

\newcolumntype{L}[1]{>{\raggedright\arraybackslash}p{#1}}

% ------------------------------------------------
% Documento
% ------------------------------------------------
\begin{document}

% ------------------------------------------------
% Cabeçalho
% ------------------------------------------------
\begin{tabularx}{\textwidth}{@{} L{4cm} X @{}}
    \raisebox{-0.5cm}{\includegraphics[width=4cm]{Universidade_de_Rio_Verde.jpg}} &
    \hfill
    \begin{tabular}{|>{\centering\arraybackslash}m{10cm}|>{\raggedright\arraybackslash}p{1.8cm}|}
        \hline
        \cellcolor{black}\color{white}\rule{0pt}{1.2cm}\textbf{LISTA DE EXERCÍCIOS: COMUNICAÇÃO VIA E-MAIL} 
        & 
        \scriptsize Nota \tabularnewline
        \hline
    \end{tabular} 
\end{tabularx}
\vspace{0.3cm}

\begin{tabularx}{\textwidth}{|X|L{2cm}|}
\hline
\textbf{Curso} & \textbf{Campus} \\
\textit{Informática Básica} & \textit{Rio Verde} \\
\hline
\multicolumn{2}{|l|}{\textbf{Disciplina}} \\  
\multicolumn{2}{|l|}{ \textit{Ferramentas de Escritório e Internet}} \\
\hline
\textbf{Nome do(a) acadêmico(a)} & \textbf{Valor} \\
& - \\
\hline
\textbf{Professor} & \textbf{Turma} \\
\textit{Sergio Souza Novak} & \textit{Matutino/Noturno} \\
\hline
\textbf{Data} & \textbf{Código} \\
\textit{13/05/2026} & \textit{INF-BAS} \\
\hline
\end{tabularx}

\vspace{0.5cm}

\begin{tcolorbox}[colback=blue!5!white,colframe=blue!75!black,title=Instruções Importantes]
Utilize sua conta do \textbf{Outlook} ou \textbf{Hotmail} para realizar as atividades abaixo. Para cada exercício, você deve preencher os campos corretamente conforme solicitado. \textbf{Atenção:} No campo "Para" (Destinatário), utilize o e-mail que o professor indicar em sala.
\end{tcolorbox}

\vspace{0.3cm}
\rule{\textwidth}{0.4pt}
\vspace{0.3cm}

% ------------------------------------------------
% EXERCÍCIO 1
% ------------------------------------------------
\section*{Exercício 1: Primeiro Contato (Nível Básico)}

\textbf{Cenário:} Você acabou de iniciar o curso e precisa enviar um e-mail de apresentação para o seu professor.

\begin{itemize}[label=$\square$]
    \item \textbf{Para:} sergio.novak@unirv.edu.br
    \item \textbf{Assunto:} Apresentação de Aluno - [Seu Nome Completo]
    \item \textbf{Corpo do E-mail:}
    \begin{tcolorbox}[colback=gray!5, colframe=gray!50]
        Olá, Professor Sergio.\\
        \\
        Gostaria de informar que estou presente na aula de Informática Básica e este é o meu e-mail oficial para contato.\\
        \\
        Atenciosamente,\\ \relax
        [Seu Nome]
    \end{tcolorbox}
\end{itemize}
\newpage
% ------------------------------------------------
% EXERCÍCIO 2
% ------------------------------------------------
\section*{Exercício 2: Dúvida Técnica com Cópia (Nível Básico)}

\textbf{Cenário:} Você está com uma dúvida sobre atalhos do Windows e quer perguntar ao professor, enviando uma cópia (CC) para um colega de classe para que ele também veja a resposta.

\begin{itemize}[label=$\square$]
    \item \textbf{Para:} sergio.novak@unirv.edu.br
    \item \textbf{CC:} [E-mail de um colega]
    \item \textbf{Assunto:} Dúvida sobre Atalhos de Teclado
    \item \textbf{Corpo do E-mail:}
    \begin{tcolorbox}[colback=gray!5, colframe=gray!50]
        Bom dia, Professor.\\
        \\
        Poderia me explicar novamente qual o atalho para alternar entre janelas abertas? Estou enviando em cópia para o meu colega que também ficou com essa dúvida.\\
        \\
        Obrigado!
    \end{tcolorbox}
\end{itemize}

\newpage

% ------------------------------------------------
% EXERCÍCIO 3
% ------------------------------------------------
\section*{Exercício 3: Convite Discreto com CCO (Nível Intermediário)}

\textbf{Cenário:} Você vai organizar um grupo de estudos e quer convidar vários colegas sem que um veja o e-mail do outro (Cópia Oculta - CCO).

\begin{itemize}[label=$\square$]
    \item \textbf{Para:} [Seu próprio e-mail]
    \item \textbf{CCO:} [E-mail de pelo menos 3 colegas]
    \item \textbf{Assunto:} Convite: Grupo de Estudos de Informática
    \item \textbf{Corpo do E-mail:}
    \begin{tcolorbox}[colback=gray!5, colframe=gray!50]
        Pessoal, boa tarde.\\
        \\
        Estou organizando um grupo para praticarmos o uso do Outlook na próxima terça-feira, após a aula. Quem tiver interesse, por favor, me responda individualmente.\\
        \\
        Até mais!
    \end{tcolorbox}
\end{itemize}

% ------------------------------------------------
% EXERCÍCIO 4
% ------------------------------------------------
\section*{Exercício 4: Envio de Arquivo/Currículo (Nível Intermediário)}

\textbf{Cenário:} Você encontrou uma vaga de estágio e precisa enviar seu currículo em anexo.

\begin{itemize}[label=$\square$]
    \item \textbf{Para:} rh.empresa@exemplo.com (Simulado)
    \item \textbf{Assunto:} Candidatura à Vaga de Estágio - [Seu Nome]
    \item \textbf{Anexo:} Um arquivo PDF qualquer renomeado para \texttt{Curriculo\_[SeuNome].pdf}.
    \item \textbf{Corpo do E-mail:}
    \begin{tcolorbox}[colback=gray!5, colframe=gray!50]
        Prezados, boa tarde.\\
        \\
        Segue em anexo o meu currículo para a vaga de estágio mencionada no portal de carreiras da UniRV.\\
        \\
        Fico à disposição para uma entrevista.\\
        \\
        Cordialmente,\\ \relax
        [Seu Telefone de Contato]
    \end{tcolorbox}
\end{itemize}

% ------------------------------------------------
% EXERCÍCIO 5
% ------------------------------------------------
\section*{Exercício 5: Relatório de Trabalho Profissional (Nível Avançado)}

\textbf{Cenário:} Você trabalha em um escritório e deve enviar o relatório semanal para seu chefe, com cópia para o setor administrativo e um arquivo Excel em anexo.

\begin{itemize}[label=$\square$]
    \item \textbf{Para:} chefe.departamento@unirv.edu.br
    \item \textbf{CC:} adm.setor@unirv.edu.br
    \item \textbf{Assunto:} Relatório Semanal de Atividades - Semana 19
    \item \textbf{Anexo:} Uma planilha do Excel ou arquivo do Word chamado \texttt{Relatorio\_S19.xlsx}.
    \item \textbf{Corpo do E-mail:}
    \begin{tcolorbox}[colback=gray!5, colframe=gray!50]
        Prezado Diretor,\\
        \\
        Conforme solicitado, envio em anexo a planilha detalhada com as atividades realizadas por nossa equipe durante esta semana. O setor administrativo já está em cópia para ciência dos prazos.\\
        \\
        Qualquer dúvida, estou à disposição.\\
        \\
        Atenciosamente,\\
        \textbf{[Seu Nome]} \\
        Assistente Administrativo
    \end{tcolorbox}
\end{itemize}

\newpage

% ------------------------------------------------
% EXERCÍCIO 6
% ------------------------------------------------
\section*{Exercício 6: Entrega de Projeto Final (Nível Desafio)}

\textbf{Cenário:} É o fim do semestre. Você deve entregar o projeto final para o Diretor do Curso, com cópia para o seu Professor e uma cópia oculta (CCO) para o seu próprio e-mail de backup. \textbf{Deve conter dois anexos.}

\begin{itemize}[label=$\square$]
    \item \textbf{Para:} diretor.informatica@unirv.edu.br
    \item \textbf{CC:} sergio.novak@unirv.edu.br
    \item \textbf{CCO:} [Seu e-mail pessoal]
    \item \textbf{Assunto:} ENTREGA FINAL: Projeto Integrador - Grupo 04
    \item \textbf{Anexos:} 
        1. \texttt{Projeto\_Escrito.pdf} \\
        2. \texttt{Apresentacao\_Slides.pptx}
    \item \textbf{Corpo do E-mail:}
    \begin{tcolorbox}[colback=gray!5, colframe=gray!50]
        Excelentíssimo Diretor,\\
        \\
        Através deste e-mail, realizamos a entrega formal dos documentos referentes ao Projeto Integrador do curso de Informática Básica. Seguem em anexo o relatório técnico e a apresentação em slides.\\
        \\
        O professor da disciplina está em cópia para a devida validação acadêmica.\\
        \\
        Agradecemos a oportunidade.\\
        \\
        Atenciosamente,\\
        \textbf{Equipe de Informática UniRV}
    \end{tcolorbox}
\end{itemize}

\vspace{1cm}
\begin{center}
    \textit{"A comunicação eficaz é a base para o sucesso profissional."} -- Boa sorte!
\end{center}

\end{document}
